\section{Dokumentation der Testdurchführung}


\subsection{Android Unit-Tests}
Die Tests zeigen, dass zentrale Funktionen der Anwendung korrekt arbeiten. 
Abgedeckt werden insbesondere:

\begin{itemize}
  \item Verarbeitung von Authentifizierungsdaten (JWT, Session-Handling)
  \item Speicherung von Einstellungen und Sitzungsdaten über SharedPreferences
  \item Konstruktion und Serialisierung von Data Classes und DTOs
  \item JSON-Verarbeitung (Serialisierung, Deserialisierung, Fehlerfälle)
  \item Fragen-Hierarchie und polymorphe Vererbung
  \item Zeit-, Alarm- und Locale-Verwaltung
  \item Utility- und Helper-Funktionen
\end{itemize}
Bei den vorliegenden Tests handelt es sich ausschließlich um lokale
Unit-Tests, welche im Verzeichnis \texttt{src/test/unit} ausgeführt werden.
Diese Tests sind bewusst plattformunabhängig gehalten und testen keine
direkten Android-UI-Komponenten oder Framework-Klassen.

Android-spezifische Funktionalitäten wie Activities, Fragments,
Navigation, Layouts oder reale Benachrichtigungen wurden nicht getestet,
da diese typischerweise über Instrumentation-Tests
(\texttt{src/androidTest}) abgesichert werden.
Die Entscheidung, den Fokus auf Unit-Tests zu legen, ermöglicht schnelle,
zuverlässige und reproduzierbare Tests der Kernlogik. Android-spezifische Abhängigkeiten (z.B. \texttt{Context}, \texttt{SharedPreferences}) werden mittels Mockito gemockt, um isolierte Tests ohne Android-Umgebung zu ermöglichen.
\subsubsection{AlleApiTests.kt}
Diese Testklasse enthält Unit Tests für verschiedene API-bezogene Komponenten unserer Android-Anwendung. Sie testet JWT-Verarbeitung, Session-Management und Data Classes für Authentifizierung und Protokoll-API.
Getestete Komponenten sind \textit{JWT Utils, Session, SessionStorage, Auth DTOsund Protocol DTOs}\\



\begin{tabularx}{\textwidth}{|l|X|l|}
\hline
\textbf{Testfall} & \textbf{ Beschreibung} & \textbf{Status} \\
\hline
1. JWT ohne Username & Gibt null zurück bei JWT ohne username-claim & Bestanden\\
\hline
2. Session Management & Setzt und liest JWT in Session-Klasse & Bestanden\\
\hline
3. Session null setzen & Setzt Session-JWT korrekt auf null zurück & Bestanden\\
\hline
4. SessionStorage speichern & Speichert JWT in SharedPreferences (Mock) & Bestanden\\
\hline
5. SessionStorage laden & Lädt JWT aus SharedPreferences (Mock) & Bestanden\\
\hline
6. SessionStorage löschen & Löscht JWT aus SharedPreferences (Mock) & Bestanden\\
\hline
7. Login Request & Testet alle Properties der Login-Data-Class & Bestanden\\
\hline
8. Token Response & Testet TokenResponse Data Class Konstruktor & Bestanden\\
\hline
9. Protocol Answer DTO & Testet einfache ProtocolAnswerDto Data Class & Bestanden\\
\hline
10. Protocol Submission Request & Testet komplexe Request mit Antwortliste & Bestanden\\
\hline
11. Protocol Submission Response & Testet Response Data Class mit 3 Feldern & Bestanden\\
\hline
12. Data Class equals() & Vergleicht zwei identische Data Class Objekte & Bestanden\\
\hline
13. Data Class toString() & Überprüft String-Repräsentation der Data Class & Bestanden\\
\hline
14. Leere Antwort & Testet Edge Case mit leerem Ergebnis-String & Bestanden\\
\hline
15. Sonderzeichen Passwort & Testet Passwort mit Sonderzeichen & Bestanden\\
\hline
16. Cleanup & Bereinigt Session nach jedem Test & Bestanden\\
\hline

\end{tabularx}

\subsubsection{AlleDataServiceTests.kt}
Diese Testklasse enthält Unit Tests für alle DataService-Komponenten der App. Sie testet Datenklassen, SharedPreferences-Operationen, JSON-Serialisierung und Fragen-Hierarchien.
Getestete Komponenten sind \textit{Answer.kt, DataLoader.kt, Question.kt, QuestionTypeAdapter.kt}\\
\begin{tabularx}{\textwidth}{|l|X|l|}
\hline
\textbf{Testfall} & \textbf{Beschreibung} & \textbf{Status} \\
\hline
1. Answer Konstruktor & Testet Answer-Data-Class mit id und result & Bestanden\\
\hline
2. QuestionSlider Properties & Testet alle Properties von QuestionSlider & Bestanden\\
\hline
3. QuestionTime Properties & Testet alle Properties von QuestionTime & Bestanden\\
\hline
4. QuestionDuration Properties & Testet alle Properties von QuestionDuration & Bestanden\\
\hline
5. QuestionMultipleChoice Properties & Testet Multiple-Choice-Fragen mit Auswahlmöglichkeiten & Bestanden\\
\hline
6. QuestionYesNo Properties & Testet Ja/Nein-Fragen mit skipIfNo Logik & Bestanden\\
\hline
7. QuestionYesNo Empty Skip & Testet Ja/Nein-Fragen ohne skipIfNo Liste & Bestanden\\
\hline
8. QuestionDate Properties & Testet alle Properties von QuestionDate & Bestanden\\
\hline
9. Question Vererbung & Prüft dass alle Fragetypen von Question erben & Bestanden\\
\hline
10. DataLoader Sprache speichern & Speichert Sprache in SharedPreferences & Bestanden\\
\hline
11. DataLoader Sprache laden & Lädt Sprache aus SharedPreferences & Bestanden\\
\hline
12. DataLoader Default Sprache & Lädt Standardsprache DE bei null & Bestanden\\
\hline
13. DataLoader Alarm Status speichern & Speichert Alarm-aktiviert Status & Bestanden\\
\hline
14. DataLoader Alarm Status laden & Lädt Alarm-aktiviert Status & Bestanden\\
\hline
15. DataLoader Alarm Zeit speichern & Speichert Alarmzeit als LocalTime & Bestanden\\
\hline
16. DataLoader Alarm Zeit laden & Lädt Alarmzeit aus SharedPreferences & Bestanden\\
\hline
17. DataLoader Default Alarmzeit & Lädt Standard-Alarmzeit (20:00) & Bestanden\\
\hline
18. DataLoader Farbschema speichern & Speichert Farbschema-Einstellung & Bestanden\\
\hline
19. DataLoader Farbschema laden & Lädt Farbschema aus SharedPreferences & Bestanden\\
\hline
20. DataLoader Default Farbschema & Lädt Standard-Farbschema system & Bestanden\\
\hline
21. Fragen-Datei nicht gefunden & Handhabt IOException bei fehlender JSON-Datei & Bestanden\\
\hline
22. Erster Start Default & Lädt Default-Wert für First-Time-Flag & Bestanden\\
\hline
23. FirstTime Completed setzen & Setzt First-Time-Flag auf false & Bestanden\\
\hline
24. Username speichern & Speichert Benutzername in SharedPreferences & Bestanden\\
\hline
25. Username laden & Lädt gespeicherten Benutzernamen & Bestanden\\
\hline
26. Gson Answer Serialisierung & Serialisiert Answer-Objekt zu JSON & Bestanden\\
\hline
\end{tabularx}
\begin{tabularx}{\textwidth}{|l|X|l|}
\hline
27. Gson Answer Deserialisierung & Deserialisiert JSON zu Answer-Objekt & Bestanden\\
\hline
28. Gson Question Liste & Serialisiert Liste von Answer-Objekten & Bestanden\\
\hline
29. Gson QuestionSlider Roundtrip & Vollständiger Serialisierungs/Deserialisierungs-Test & Bestanden\\
\hline
30. QuestionTypeAdapter TypeMap & Testet TypeAdapter für Frage-Typen & Bestanden\\
\hline
31. Deserialisierung QuestionSlider & Deserialisiert QuestionSlider aus JSON & Bestanden\\
\hline
32. Deserialisierung MultipleChoice & Deserialisiert QuestionMultipleChoice aus JSON & Bestanden\\
\hline
33. Deserialisierung QuestionYesNo & Deserialisiert QuestionYesNo aus JSON & Bestanden\\
\hline
34. Deserialisierung QuestionTime & Deserialisiert QuestionTime aus JSON & Bestanden\\
\hline
35. Deserialisierung QuestionDuration & Deserialisiert QuestionDuration aus JSON & Bestanden\\
\hline
36. Deserialisierung QuestionDate & Deserialisiert QuestionDate aus JSON & Bestanden\\
\hline
37. Deserialisierung ungültiger Typ & Wirft Exception bei unbekanntem Typ & Bestanden\\
\hline
38. Deserialisierung ohne Typ & Wirft Exception bei fehlendem Typ-Feld & Bestanden\\
\hline
39. Deserialisierung leerer Typ & Wirft Exception bei leerem Typ-Feld & Bestanden\\
\hline
40. DataLoader Antworten speichern & Schreibt Antworten in Datei & Bestanden\\
\hline
41. Frage-Hierarchie & Testet Vererbungshierarchie aller Fragetypen & Bestanden\\
\hline
42. Answer Listen Operationen & Testet Listen-Operationen mit Answer-Objekten & Bestanden\\
\hline
43. Edge Cases - leere Strings & Testet mit leeren Strings in Antworten & Bestanden\\
\hline
44. Edge Cases - Sonderzeichen & Testet Umlaute und Sonderzeichen in Texten & Bestanden\\
\hline
45. Edge Cases - Maximalwerte & Testet mit Maximalwerten für IDs und Floats & Bestanden\\
\hline
46. Edge Cases - negative Werte & Testet mit negativen Werten für Slider & Bestanden\\
\hline
47. Liste contains Operation & Testet Listen-Operationen in QuestionYesNo & Bestanden\\
\hline
48. JSON Format mit Quotes & Testet JSON-Serialisierung mit Anführungszeichen & Bestanden\\
\hline
\end{tabularx}



\subsubsection{AlleHelperTests.kt}
Diese Testklasse enthält Unit Tests für alle Helper-Klassen der App. Sie testet Notification-Logik, Alarm-Berechnungen, Locale-Verwaltung und verschiedene Hilfsfunktionen.
Getestete Komponenten sind \textit{NotificationHelper, AlarmReceiver (indirekt), updateLocale Funktion und Validierungs-Helfer}\\
\vspace{1em}

\begin{tabularx}{\textwidth}{|l|X|l|}
\hline
\textbf{Testfall} & \textbf{Beschreibung} & \textbf{Status} \\
\hline
1. Alarm-Berechnung heute & Berechnet Alarmzeit für heute (Zeit in Zukunft) & Bestanden\\
\hline
2. Alarm-Berechnung morgen & Berechnet Alarmzeit für morgen (Zeit bereits vorbei) & Bestanden\\
\hline
3. Millisekunden Konvertierung & Konvertiert LocalDateTime zu Millisekunden & Bestanden\\
\hline
4. Alarm Edge Cases & Testet Mitternacht und gleiche Zeiten & Bestanden\\
\hline
5. AlarmReceiver Konstanten & Testet Notification-ID (hardcoded) & Bestanden\\
\hline
6. Notification Properties & Testet Channel-ID und Priority & Bestanden\\
\hline
7. updateLocale Deutsch & Testet deutsche Locale Erstellung & Bestanden\\
\hline
8. updateLocale Englisch & Testet englische Locale Erstellung & Bestanden\\
\hline
9. updateLocale Default & Testet Default Locale & Bestanden\\
\hline
10. Zeitvalidierung gültig & Testet gültige Uhrzeiten (00:00 - 23:59) & Bestanden\\
\hline
11. Zeitvalidierung ungültig & Testet ungültige Uhrzeiten (negative/zu große Werte) & Bestanden\\
\hline
12. Zeitformatierung & Formatiert Uhrzeiten als HH:mm und mit "Uhr" & Bestanden\\
\hline
13. Datumsformatierung & Formatiert Datum als TT.MM.JJJJ & Bestanden\\
\hline
14. Intent Actions & Testet Android Intent-Action Konstanten & Bestanden\\
\hline
15. Package URI Format & Testet URI Format für Package-Namen & Bestanden\\
\hline
16. SharedPreferences Keys & Testet Konsistenz der SharedPreferences Keys & Bestanden\\
\hline
17. Mock Context & Testet Mock-Setup für Context und SharedPreferences & Bestanden\\
\hline
18. Android Version Checks & Testet Android SDK Version Konstanten & Bestanden\\
\hline
19. Alarm-Berechnung Helper & Testet Hilfsfunktion für Alarmzeit-Berechnung & Bestanden\\
\hline
20. Locale Funktionalität & Testet Locale-Klassenmethoden & Bestanden\\
\hline
21. String Resources Mock & Testet String-Resource-Namen & Bestanden\\
\hline
22. Edge Cases Zeiten & Testet Extremfälle für Uhrzeiten & Bestanden\\
\hline
23. Edge Cases Daten & Testet besondere Daten (Schaltjahre, 2038-Problem) & Bestanden\\
\hline
24. TimeZone Handling & Testet Zeitzonen-Verarbeitung & Bestanden\\
\hline
25. PendingIntent Flags & Testet Kombination von PendingIntent Flags & Bestanden\\
\hline
\end{tabularx}
\subsection{Android Integrationstests}
Die durchgeführten Integrationstests folgen einer use-case-basierten Teststrategie, 
bei der typische Nutzungsszenarien der App mit der Backend-API simuliert werden. 
Jeder Test repräsentiert einen konkreten Anwendungsfall aus der Benutzerperspektive.

\noindent\textbf{Technische Architektur:}

\begin{itemize}
  \item \textbf{MockWebServer:} Lokaler HTTP-Server von OkHttp, der API-Antworten simuliert
  \item \textbf{Retrofit:} HTTP-Client für API-Kommunikation
  \item \textbf{Gson:} JSON-Serialisierung/Deserialisierung
  \item \textbf{Kotlin Coroutines:} Asynchrone Ausführung der API-Aufrufe
  \item \textbf{Test-Location:} \texttt{src/test/integration/}
\end{itemize}

\noindent\textbf{Test-Strategie:}

\begin{itemize}
  \item \textbf{Success-Fälle:} Korrekte Requests führen zu erwarteten Responses
  \item \textbf{Error-Fälle:} HTTP-Fehler (400, 401, 404) werden korrekt behandelt
  \item \textbf{Validierung:} Request-Methoden, Pfade, Headers und Body-Inhalte
  \item \textbf{Edge-Cases:} Ungültige JSON-Strukturen, fehlende Felder, leere Antworten
\end{itemize}

\noindent\textbf{Setup-Prozess:}
\begin{enumerate}
  \item MockWebServer starten auf zufälligem Port
  \item Retrofit-Instanz mit MockWebServer-URL konfigurieren
  \item API-Interface über Retrofit instanziieren
  \item Mock-Antworten mit \texttt{MockResponse} vorbereiten
  \item Tests mit \texttt{runBlocking} für Coroutine-Support ausführen
  \item MockWebServer nach jedem Test herunterfahren
\end{enumerate}

\subsubsection{AuthApiIntegrationTest.kt}
Testet die Authentifizierungs-API-Endpunkte für Patient Login und First-Login.

\textbf{Getestete Endpunkte:}
\begin{itemize}
  \item \texttt{POST /api/auth/patient/first-login} - Erstanmeldung mit Aktivierungscode
  \item \texttt{POST /api/auth/patient/login} - Reguläre Anmeldung mit Username/Password
\end{itemize}

\textbf{Test-Szenarien:}
\begin{itemize}
  \item Erfolgreiche Authentifizierung mit korrektem Token
  \item Fehlerhafte Anfragen (ungültiger Code, falsche Credentials)
  \item Server-Fehler (400 Bad Request, 401 Unauthorized)
  \item Ungültige JSON-Responses
  \item Sequenzielle API-Aufrufe
\end{itemize}

\begin{tabularx}{\textwidth}{|l|X|l|}
\hline
\textbf{Testfall} & \textbf{Beschreibung} & \textbf{Status} \\
\hline
1. Patient First-Login Success & Erfolgreiche Erstanmeldung mit Token-Rückgabe & Bestanden\\
\hline
2. Patient Login Success & Erfolgreiche reguläre Anmeldung & Bestanden\\
\hline
3. First-Login Invalid Code & HTTP 400 bei ungültigem Aktivierungscode & Bestanden\\
\hline
4. Login Wrong Credentials & HTTP 401 bei falschen Anmeldedaten & Bestanden\\
\hline
5. Invalid JSON Response & Exception bei ungültiger JSON-Struktur & Bestanden\\
\hline
6. Multiple Sequential Requests & Korrekte Verarbeitung sequentieller Aufrufe & Bestanden\\
\hline
\end{tabularx}

\subsubsection{ProtocolApiIntegrationTest.kt}
Testet das Absenden von ausgefüllten Schlafprotokollen an die Backend-API.

\textbf{Getesteter Endpunkt:}
\begin{itemize}
  \item \texttt{POST /api/protocols/submissions} - Übermittlung ausgefüllter Protokolle
\end{itemize}

\textbf{Test-Szenarien:}
\begin{itemize}
  \item Erfolgreiche Protokoll-Übermittlung mit Submission-ID
  \item Ungültige Protokolldaten (HTTP 400)
  \item Fehlende/ungültige Authentifizierung (HTTP 401)
  \item Validierung von Request-Body und JSON-Struktur
\end{itemize}

\begin{tabularx}{\textwidth}{|l|X|l|}
\hline
\textbf{Testfall} & \textbf{Beschreibung} & \textbf{Status} \\
\hline
1. Submit Protocol Success & Erfolgreiche Übermittlung mit Submission-ID & Bestanden\\
\hline
2. Submit Invalid Data & HTTP 400 bei ungültigen Protokolldaten & Bestanden\\
\hline
3. Submit Unauthorized & HTTP 401 bei fehlender Authentifizierung & Bestanden\\
\hline
\end{tabularx}

\subsubsection{ReportApiIntegrationTest.kt}
Testet den Download von PDF-Berichten mit verschiedenen Parametern.

\textbf{Getesteter Endpunkt:}
\begin{itemize}
  \item \texttt{GET /api/reports/me/pdf} - Download persönlicher PDF-Berichte
\end{itemize}

\textbf{Test-Szenarien:}
\begin{itemize}
  \item Erfolgreicher PDF-Download mit korrekten Headers
  \item Verschiedene Query-Parameter (Date-Range, days, allSubmissions)
  \item Fehlende Authentifizierung (HTTP 401)
  \item Keine Daten verfügbar (HTTP 404)
  \item Validierung von Content-Type und Content-Disposition Headers
\end{itemize}

\begin{tabularx}{\textwidth}{|l|X|l|}
\hline
\textbf{Testfall} & \textbf{Beschreibung} & \textbf{Status} \\
\hline
1. PDF Download Default & Erfolgreicher Download mit Standard-Parametern & Bestanden\\
\hline
2. PDF with Date Range & Download mit spezifischem Zeitraum & Bestanden\\
\hline
3. PDF Unauthorized & HTTP 401 bei fehlender Authentifizierung & Bestanden\\
\hline
4. PDF No Data & HTTP 404 wenn keine Daten vorhanden &Bestanden\\
\hline
\end{tabularx}
\subsection{Backend-Modultests des Schlaftagebuch-Systems (Spring Boot)}

\subsubsection{Getestete Module und Use Cases}

Die Tests befinden sich unter \texttt{test/java/de.schlaftagebuch.schlaftagebuch\_backend/[Modul]} und umfassen die vier Kernmodule des Systems:

\begin{table}[h]
\centering
\begin{tabular}{|l|p{0.60\textwidth}|}
\hline
\textbf{Modul} & \textbf{Funktionalität} \\
\hline
\textbf{Auth} & Arzt-/Patient-Login, Registrierung, JWT-Token, Passwort-Verschlüsselung \\
\hline
\textbf{ArztManagement} & Code-Generierung für Patienten, Arzt-Patienten-Zuordnung, Code-Verwaltung \\
\hline
\textbf{Protocol} & Erfassung und Verwaltung von Schlaftagebuch-Einträgen (CRUD) \\
\hline
\textbf{Report} & PDF-Generierung, Datenanalyse, Visualisierung, Export \\
\hline
\end{tabular}
\caption{Getestete Module entsprechend der Use Cases}
\end{table}

Diese Module bilden die in den Use Cases definierten Hauptfunktionalitäten ab und enthalten die kritischen Sicherheits- und Geschäftslogik-Komponenten.

\subsubsection{Teststrategie}

Die Testimplementierung basiert auf einer \textbf{rein Unit-Test-basierten Strategie} mit Mockito. Services und Controller werden isoliert getestet, wobei alle Abhängigkeiten (Repositories, externe Services) gemockt werden. Diese fokussierte Herangehensweise ermöglicht schnelle Testausführungen und konzentriert sich auf die Validierung der Business-Logik unter kontrollierten Bedingungen.
\subsubsection{Testszenarien}

Die gewählten Unit-Test-Szenarien decken die Kernfunktionalitäten umfassend ab:

\begin{itemize}
    \item \textbf{Positivfälle:} Testung aller erwarteten Erfolgsszenarien mit gültigen Eingabedaten
    \item \textbf{Negativfälle:} Systematische Testung von Grenzwerten und ungültigen Eingaben (24\% der Tests)
    \item \textbf{Edge Cases:} Behandlung von Sonderfällen wie Null-Werten, leeren Listen, extremen Datenmengen
    \item \textbf{Fehlerbehandlung:} Verifikation der korrekten Exception-Propagation von Service- zu Controller-Schicht
    \item \textbf{Verhaltensverifikation:} Überprüfung der korrekten Interaktion mit gemockten Abhängigkeiten
\end{itemize}

\subsubsection{Technische Infrastruktur}
\begin{itemize}
    \item Spring Boot 3.4.0, Java 17
    \item \textbf{JUnit 5 + Mockito 5.x} für reine Unit-Tests
    \item AssertJ für flüssige Assertions
    \item Gradle mit JaCoCo Plugin für Coverage-Analyse
\end{itemize}

\subsubsection{Ausführung}

\begin{verbatim}
# Alle Tests
./gradlew clean test

# Mit Coverage-Report
./gradlew test jacocoTestReport
\end{verbatim}

Report: \texttt{build/reports/jacoco/test/html/index.html}
\subsubsection{Testergebnisse}

\begin{table}[H]
\centering
\begin{tabular}{|l|c|}
\hline
\textbf{Getesteter Funktionsbereich} & \textbf{Status} \\
\hline
Arzt-Login (Erfolg, falsche Daten, falsche Rolle) & \checkmark \\
\hline
Patienten-Login (Erfolg, falsche Daten, falsche Rolle) & \checkmark \\
\hline
Einmal-Code-Erzeugung (neu / bestehender Patient) & \checkmark \\
\hline
Patienten-Erstlogin mit Einmal-Code & \checkmark \\
\hline
Validierung von Codes (ungültig, abgelaufen, benutzt, widerrufen) & \checkmark \\
\hline
Validierung von Eingabedaten (Geburtsdatum, Consent) & \checkmark \\
\hline
JWT-Token-Erstellung und Rollenprüfung & \checkmark \\
\hline
AuthController-Endpunkte (Service-Delegation, Fehlerweitergabe) & \checkmark \\
\hline
\end{tabular}
\caption{Testergebnisse des Auth-Moduls}
\end{table}



\begin{table}[H]
\centering
\begin{tabular}{|l|c|}
\hline
Anzahl Unit-Tests (Service) & 21\\
Anzahl Unit-Tests (Controller) & 5\\
Bestandene Tests & 26(100\%) \\
Fehlgeschlagene Tests & 0 \\
Code Coverage (Service-Schicht) & 75\% \\
\hline
\end{tabular}
\caption{Zusammenfassung der Testergebnisse für das Auth-Modul}
\end{table}


\begin{table}[H]
\centering
\begin{tabular}{|l|c|}
\hline
\textbf{Getesteter Funktionsbereich} & \textbf{Status} \\
\hline
Verwaltung von Patienten- und Arzt-Zuordnungen & \checkmark \\
\hline
Behandlung fehlender oder ungültiger Patientendaten & \checkmark \\
\hline
Generierung, Statusverwaltung und Deaktivierung von Einmal-Codes & \checkmark \\
\hline
Berechnung und Rückgabe von Dashboard-Statistiken & \checkmark \\
\hline
Delegation der REST-Endpunkte an die jeweilige Service-Schicht & \checkmark \\
\hline

\end{tabular}
\caption{Testergebnisse des ArztManagement-Moduls}
\end{table}

\begin{table}[H]
\centering
\begin{tabular}{|l|c|}
\hline
\textbf{Metrik} & \textbf{Wert} \\
\hline
Anzahl Unit-Tests (ArztDashboardService) & 10 \\
Anzahl Unit-Tests (ArztManagementService) & 20 \\
Anzahl Unit-Tests (ArztManagementController) & 8\\
Bestandene Tests & 38 (100\%) \\
Fehlgeschlagene Tests & 0 \\
Code Coverage (Service-Schicht) & 87\% \\
\hline
\end{tabular}
\caption{Zusammenfassung der Testergebnisse für das ArztManagement-Modul}
\end{table}

\begin{table}[H]
\centering
\begin{tabular}{|l|c|}
\hline
\textbf{Getesteter Funktionsbereich} & \textbf{Status} \\
\hline
Verarbeitung von Protokoll-Einreichungen verschiedener Benutzertypen & \checkmark \\
\hline
Validierung gültiger Anfragen & \checkmark \\
\hline
JSON-Serialisierung mit Sonderzeichen und Escaping & \checkmark \\
\hline
Unterstützung verschiedener Protokoll-Vorlagen und Sprachen & \checkmark \\
\hline
Automatische Zeitstempel-Generierung bei fehlenden Daten & \checkmark \\
\hline
Behandlung von Edge Cases und großen Datenmengen & \checkmark \\
\hline
Korrektes Mapping zwischen Entity und Response-DTO & \checkmark \\
\hline
Exception-Propagation von Service zu Controller & \checkmark \\
\hline
\end{tabular}
\caption{Testergebnisse des Protocol-Moduls}
\end{table}

Im Rahmen der Testentwicklung wurden zusätzlich Sonderfälle mit
\texttt{null}- oder leeren Eingaben berücksichtigt, um die Vollständigkeit
der Testbetrachtung sicherzustellen.

\begin{itemize}
    \item Übergabe eines \texttt{null}-Request-Objekts
    \item Übergabe einer \texttt{null}-Antwortliste
    \item Übergabe einer leeren Antwortliste
\end{itemize}

Diese Szenarien treten im produktiven Betrieb nicht auf, da die Request-Objekte
in der Controller-Schicht durch Spring Boot erzeugt, validiert und mit
Default-Werten initialisiert werden.

\textbf{Hinweis:} Die genannten Sonderfälle dienen der theoretischen
Robustheitsbetrachtung und haben keine Auswirkung auf den realen Betrieb
des Systems.

\begin{table}[H]
\centering
\begin{tabular}{|l|c|}
\hline
\textbf{Metrik} & \textbf{Wert} \\
\hline
Anzahl Unit-Tests (Service) & 21 \\
Anzahl Unit-Tests (Controller) & 8 \\
Gesamtzahl Unit-Tests & 29 \\
Bestandene Tests & 26 (90\%) \\
Fehlgeschlagene Tests & 3 \\
Code Coverage (Service-Schicht) & 94\% \\
Code Coverage (Controller-Schicht) & 89\% \\
\hline
\end{tabular}
\caption{Zusammenfassung der Testergebnisse für das Protocol-Modul}
\end{table}

\begin{table}[H]
\centering
\begin{tabular}{|l|c|}
\hline
\textbf{Getesteter Funktionsbereich} & \textbf{Status} \\
\hline
Report-Generierung für Patienten (eigene Berichte) & \checkmark \\
\hline
Report-Generierung für Ärzte (Patientenberichte) & \checkmark \\
\hline
PDF-Export von Schlaftagebuch-Berichten & \checkmark \\
\hline
Zeitraum-basierte Report-Generierung (Date Range) & \checkmark \\
\hline
Tages-basierte Report-Generierung (Last N Days) & \checkmark \\
\hline
Automatische Patientenerkennung aus Authentifizierung & \checkmark \\
\hline
Behandlung verschiedener Eingabeformate (JSON) & \checkmark \\
\hline
Validierung von Eingabeparametern & \checkmark \\
\hline
Berechnung von Schlafstatistiken & \checkmark \\
\hline
Integration mit FHIR für medizinische Berichte & \checkmark \\
\hline
\end{tabular}
\caption{Testergebnisse des Report-Moduls}
\end{table}


\begin{table}[H]
\centering
\begin{tabular}{|l|c|}
\hline
\textbf{Metrik} & \textbf{Wert} \\
\hline
Anzahl Unit-Tests (Service) & 23 \\
Anzahl Unit-Tests (Controller) & 9 \\
Gesamtzahl Unit-Tests & 32 \\
Bestandene Tests & 32 (100\%) \\
Fehlgeschlagene Tests & 0\\
Code Coverage (Service-Schicht) & 94\% \\
Code Coverage (Controller-Schicht) & 89\% \\
\hline
\end{tabular}
\caption{Zusammenfassung der Testergebnisse für das Report-Modul}
\end{table}



\subsection{Test mit Postman}

Im Folgenden wird die Testdurchführung der relevanten API-Endpunkte mit Postman dokumentiert. Die Tests erfolgen lokal gegen \texttt{http://localhost:8080}. Für alle Requests mit JWT ist sicherzustellen, dass das Token vollständig kopiert und als Bearer Token gesetzt wird.

\subsubsection{Als Arzt einloggen}
\textbf{Methode:} \texttt{POST}\\
\textbf{URL:} \texttt{http://localhost:8080/api/auth/arzt/login}

\textbf{Body (JSON):}
\begin{verbatim}
{
  "username": "arzt1",
  "password": "arztpass"
}
\end{verbatim}

\textbf{Erwartete Antwort:} Es wird ein JWT zurückgegeben, z.\,B.:
\begin{verbatim}
{ "token": "eyJhbGciOi..." }
\end{verbatim}
Der Token muss vollständig kopiert werden.

\subsubsection{Zugangscode generieren}
\textbf{Methode:} \texttt{POST}\\
\textbf{URL:} \texttt{http://localhost:8080/api/auth/patients/generate-code}

\textbf{Auth in Postman:} \texttt{Auth} $\rightarrow$ \texttt{Bearer Token} und das zuvor erhaltene Arzt-JWT einfügen.\\
\textbf{Body:} kein Body

\textbf{Erwartete Antwort:} \texttt{GenerateCodeResponse} mit \texttt{code} (und ggf. \texttt{patientId}).

\subsubsection{Erster Login des Patienten}
\textbf{Methode:} \texttt{POST}\\
\textbf{URL:} \texttt{http://localhost:8080/api/auth/patient/first-login}

\textbf{Body (JSON, Code ersetzen):}
\begin{verbatim}
{
  "code": "MQ4842",
  "consentAccepted": true,
  "givenName": "Test",
  "familyName": "Patient",
  "password": "pass123",
  "birthDate": "2000-01-01"
}
\end{verbatim}

\textbf{Ergebnis:} Der Patient wird initial registriert/aktiviert (abhängig von Backend-Implementierung) und kann anschließend authentifiziert werden.

\subsubsection{Test-Login Patient (über Seed, DEV)}
\textbf{Hinweis:} Funktioniert aktuell \emph{ohne} Authentifizierung und dient ausschließlich Testzwecken in der Entwicklungsumgebung.\\
\textbf{Methode:} \texttt{POST}\\
\textbf{URL:} \texttt{http://localhost:8080/api/auth/patient/login}

\textbf{Body (JSON):}
\begin{verbatim}
{
  "username": "patient_test",
  "password": "patientpass"
}
\end{verbatim}

\textbf{Erwartete Antwort:} Enthält das JWT des Patienten (Patient-Token), welches für die folgenden geschützten Endpunkte verwendet wird.

\subsubsection{Protokoll absenden}
\textbf{Methode:} \texttt{POST}\\
\textbf{URL:} \texttt{http://localhost:8080/api/protocols/submissions}

\textbf{Authorization Header:} \texttt{Authorization: Bearer <PATIENT\_TOKEN>}

\textbf{Body (JSON):}
\begin{verbatim}
{
  "templateKey": "schlaftagebuch_v1",
  "locale": "de",
  "filledAt": "2025-12-30T10:15:00Z",
  "answers": [
    { "id": 0, "result": "7" },
    { "id": 1, "result": "22:45" }
  ]
}
\end{verbatim}

\textbf{Erwartetes Ergebnis:} Wenn eine \texttt{submissionId} zurückkommt, wurde das Protokoll gespeichert und dem Patienten zugeordnet.

\subsubsection{Bericht erstellen (PDF) -- Test}
\textbf{Methode:} \texttt{GET}\\
\textbf{URL:} \texttt{http://localhost:8080/api/reports/me/pdf?days=14}

\textbf{Authorization Header:} \texttt{Authorization: Bearer <PATIENT\_TOKEN>}

\textbf{Ergebnis:} Es wird ein PDF-Bericht für den Patienten generiert (Zeitraum: letzte \texttt{days=14} Tage), sofern entsprechende Daten vorhanden sind.

\newpage

\subsection{Query-Parameter}

\textbf{Warum benötigen wir einen Query-Parameter?}\\
Im Produktivbetrieb ist es sinnvoll, im PDF-Bericht pro Kalendertag ausschließlich
den jeweils neuesten Schlafprotokoll-Eintrag zu berücksichtigen. Dadurch bleibt
der Bericht übersichtlich, gut vergleichbar und medizinisch sinnvoll interpretierbar.

Für Test- und Entwicklungszwecke ist dieses Verhalten jedoch unpraktisch, da man
ansonsten viele Kalendertage abwarten müsste, um ausreichend Daten für die
Berichtserstellung und Visualisierung zu erzeugen.

Aus diesem Grund wird ein zusätzlicher \textbf{Query-Parameter} eingeführt, mit dem
das Verhalten der Berichtserstellung gezielt zwischen Produktiv- und Testmodus
umgeschaltet werden kann.

\begin{itemize}
  \item \texttt{allSubmissions=false} (Standardbetrieb):\\
        Pro Kalendertag wird nur die neueste Submission im Bericht berücksichtigt.
        Dieses Verhalten ist für den regulären Einsatz vorgesehen.
  \item \texttt{allSubmissions=true} (Testmodus):\\
        Alle Submissions im gewählten Zeitraum werden berücksichtigt, auch wenn
        mehrere Einträge am selben Tag existieren.
\end{itemize}

\textbf{Wie funktioniert das technisch?}\\
Ein Query-Parameter ist ein optionaler Parameter, der bei einem HTTP-Request an den
Endpoint angehängt wird. Er beginnt mit einem Fragezeichen (\texttt{?}); mehrere
Parameter werden mit einem kaufmännischen Und-Zeichen (\texttt{\&}) getrennt.

Im Backend wird der Parameter über das Spring-Framework mit
\texttt{@RequestParam} entgegengenommen und steuert dort die interne Logik der
Datenaggregation.

\paragraph{Relevanter Endpoint}
Der PDF-Bericht wird über folgenden Endpoint erzeugt:
\[
\texttt{GET /api/reports/me/pdf}
\]

\paragraph{Beispiele für die Verwendung}
\begin{itemize}
  \item \textbf{Standardbetrieb (Produktivmodus, 14 Tage):}
  \[
  \texttt{/api/reports/me/pdf?days=14}
  \]

  \item \textbf{Testmodus (alle Submissions im Zeitraum):}
  \[
  \texttt{/api/reports/me/pdf?days=14\&allSubmissions=true}
  \]

  \item \textbf{Schnelltest (ein Tag, mehrere Submissions):}
  \[
  \texttt{/api/reports/me/pdf?days=1\&allSubmissions=true}
  \]
\end{itemize}

\textbf{Definition im Backend}\\
Der Parameter ist im Controller wie folgt definiert:
\[
\texttt{@RequestParam(defaultValue = "false") boolean allSubmissions}
\]

Dadurch ist sichergestellt, dass ohne explizite Angabe automatisch das
Standardverhalten (Produktivmodus) verwendet wird.

\paragraph{Konkretes Testvorgehen}
\begin{enumerate}
  \item \textbf{Erzeugen mehrerer Submissions:}\\
  Das Schlafprotokoll innerhalb der mobilen Anwendung mehrfach am selben Tag
  ausfüllen und absenden.

  \item \textbf{Abrufen des PDF-Berichts im Testmodus:}\\
  Der Bericht wird innerhalb der App mit aktivierter Testkonfiguration
  (\texttt{allSubmissions=true}) generiert.

  \item \textbf{Erwartetes Ergebnis:}\\
  Im generierten PDF erscheinen mehrere Einträge für denselben Kalendertag.
  Die Diagramme enthalten entsprechend mehrere Datenpunkte mit eindeutig
  unterscheidbaren Beschriftungen.

  \item \textbf{Rückkehr zum Standardbetrieb:}\\
  Wird der Parameter \texttt{allSubmissions} nicht gesetzt oder auf
  \texttt{false} belassen, enthält der Bericht wieder ausschließlich den
  jeweils neuesten Eintrag pro Tag.
\end{enumerate}












